\section*{Chapter \ref{ch:m1:getting-access-futures-options} --- Getting Access: Futures and Listed Options}

\begin{solution}{exo:m1:getting-access-futures-options:1}
Non-member: $10\,000 \times 1.45 = \$14\,500$. Lessee: $10\,000 \times 0.57 + 1\,500
= \$7\,200$.
\end{solution}

\begin{solution}{exo:m1:getting-access-futures-options:2}
Lease against outside: $1\,500/(1.45 - 0.57) = 1\,705$ sides a month. Owning
against leasing: $(4\,000 - 1\,500)/(0.57 - 0.45) = 20\,833$ sides a month.
\end{solution}

\begin{solution}{exo:m1:getting-access-futures-options:3}
E-mini: $2 \times 1.20/12.50 = 19\%$; $2 \times 0.47/12.50 = 8\%$; $2 \times
0.35/12.50 = 6\%$. Micro: $2 \times 0.22/1.25 = 35\%$; $2 \times 0.07/1.25 = 11\%$;
$2 \times 0.04/1.25 = 6\%$.
\end{solution}

\begin{solution}{exo:m1:getting-access-futures-options:4}
8\,400 round trips earn $8\,400 \times 0.3 \times 12.50 = \$31\,500$ gross.
Non-member: fees $16\,800 \times 1.45 = \$24\,360$: net \$7\,140. Lessee: $16\,800
\times 0.57 + 1\,500 = \$11\,076$: net \$20\,424. The same strategy is almost
three times as profitable.
\end{solution}

\begin{solution}{exo:m1:getting-access-futures-options:5}
At the normal pace it ends the month at 49\,000, a thousand short. Reaching
50\,000 earns 5 cents on all 50\,000 sides: \$2\,500, or \$2.50 for each of the
last 1\,000 sides, several times the fee on them. It should trade them if it
can do so at a loss of less than that.
\end{solution}

\begin{solution}{exo:m1:getting-access-futures-options:6}
For large intraday positions the second is dearer in the way that matters:
ten minutes to meet a call means holding the cash idle in advance, or being
liquidated; two hours allows funding from elsewhere. The 150\% applies to
overnight positions, which are small. Ask both: how intraday exposure is
measured and limited; what triggers liquidation and who decides; interest
paid on excess collateral; which collateral is accepted; and how they
behaved towards clients in the last stressed month.
\end{solution}

\begin{solution}{exo:m1:getting-access-futures-options:7}
The outside route now has a fixed cost of \$1\,200: the break-even falls to
$(1\,500 - 1\,200)/0.88 = 341$ sides a month. With a lease at \$1\,100 the lease
is cheaper at any volume: its fixed cost is lower \emph{and} its variable
cost is lower, and the function returns zero.
\end{solution}

\begin{solution}{exo:m1:getting-access-futures-options:8}
Five sixths of the profit is a stipend. Ask: for how long is the programme
guaranteed, and can the exchange change or end it; is the stipend per firm
or does it scale with volume (tenfold capital does not earn tenfold a fixed
stipend); how many firms share it; what are the obligations and what do they
cost in a fast market; what is the \$30\,000 of spread income after the
\omterm{def:m1:what-a-trading-firm-does:adverse-selection}{adverse selection} that obligations impose. The scalable business here is the
\$30\,000, and it may be negative without the programme.
\end{solution}

\begin{solution}{pb:m1:getting-access-futures-options:1}
\textbf{1.} 12\,600 sides; $12\,600 \times 0.25 \times 12.50 = \$39\,375$.
\textbf{2.} $12\,600 \times 1.45 = \$18\,270$.
\textbf{3.} Net \$21\,105; fees are 46\% of gross.
\textbf{4.} $2 \times 1.45/12.50 = 0.23$ tick a round trip, against a gross of
0.50.
\textbf{5.} 1\,705 sides a month, 81 a day: the firm trades seven times that.
\textbf{6.} $12\,600 \times 0.57 + 1\,500 = \$8\,682$.
\textbf{7.} Net \$30\,693: \$9\,588 a month better, 45\% more.
\textbf{8.} $9\,588/21 = \$457$ a day: the application fee is repaid in 4.4
trading days.
\textbf{9.} An application with disclosure of background and finances, the
exchange's rules and disciplinary jurisdiction applying to the member
personally, restrictions on whose account may be traded at \omterm{def:m1:getting-access-futures-options:membership}{member rates} (the
member's own), and the administrative link between the lease and the
account at the FCM.
\textbf{10.} 20\,833 sides a month, about 990 a day.
\textbf{11.} 42\,000 sides; gross \$13\,125. Non-member: $42\,000 \times 0.32 =
\$13\,440$: net $-\$315$. Lessee: $42\,000 \times 0.12 = \$5\,040$: net
$+\$8\,085$ (the lease is already paid for by the E-mini).
\textbf{12.} That in small contracts, whose fees are a third of a tick for
outsiders, short-horizon strategies exist only for members. A non-member
must hold positions long enough for fees to be negligible against the move.
\textbf{13.} Nothing: the firm trades 12\,600 sides, under the threshold of
20\,000. It would need 952 sides a day.
\textbf{14.} As zero. Income that depends on a programme the exchange can end
is a bonus, not a basis for hiring or for capital commitments.
\textbf{15.} Price discrimination in favour of those who bring volume and
liquidity, and a legacy of the exchanges' origin as member-owned clubs:
memberships were the ownership. The non-member rate is what the customer of
convenience pays.
\textbf{16.} Former floor traders, their estates, firms that bought more
seats than they use, and investors: a seat yields a lease income and may
carry other rights. Lessees are firms whose volume justifies the rate but
not the purchase price.
\textbf{17.} The margin. At 600 sides a day the lease is worth \$9\,588 a
month; a jump from 120\% to 200\% of exchange margin may halve the size the
firm can carry, which halves the gross of \$39\,375, or forces it to raise
capital.
\textbf{18.} Volume large enough that FCM commissions and margin multiples
cost more than clearing-house membership: capital requirements, default-fund
contributions, operations staff, and direct exposure to other members'
defaults. For a two-person firm, never; for a large proprietary firm, a
standard step.
\textbf{19.} \textbf{1\,705 sides a month.}
\textbf{20.} A lower price per trade, in exchange for a fixed cost and for
accepting the exchange's jurisdiction.
\end{solution}

\begin{solution}{iq:m1:getting-access-futures-options:1}
A \omterm{def:m1:getting-access-futures-options:fcm}{futures commission merchant} accepts customers' orders and their money. For
a trading firm it is the \omterm{def:m1:getting-access-futures-options:fcm}{clearing member} that carries its positions at the
clearing house, collects and posts margin, guarantees its trades, provides
or accepts executions by \omterm{def:m1:getting-access-equities:brokers}{give-up}, and sets the firm's real limits: margin
multiple, intraday exposure, liquidation rights.
\iqlookfor{the FCM as credit provider, not as order router.}
\end{solution}

\begin{solution}{iq:m1:getting-access-futures-options:2}
Exchanges were member-owned; membership rates survive as a volume discount
attached to a tradable or leasable seat. The gap is large: for the main
equity-index contract a non-member pays roughly three times an owner's
exchange fee per side, and the gap is wider still relative to the tick in
the smallest contracts.
\iqlookfor{an order of magnitude for the gap, and the lease as the way in.}
\end{solution}

\begin{solution}{iq:m1:getting-access-futures-options:3}
Fees in ticks at the rate the firm will actually pay: for a non-member a
micro round trip costs a third of a tick in exchange fees alone, plus
commission. Then queue position (a one-tick scalp needs fills at the front of
the queue, which the backtest probably assumed), latency, and whether the
fills assumed are achievable at the touch at all. Most such strategies die at
the first check.
\iqlookfor{fees as a fraction of the tick, by status.}
\end{solution}

\begin{solution}{iq:m1:getting-access-futures-options:4}
Lease when monthly volume times the per-side saving exceeds the lease,
usually a few dozen round trips a day in the main contracts. Buy when the
saving of owner over lessee rates, times volume, exceeds the cost of the
capital tied up net of the lease paid, and the firm expects to stay in the
product for years: memberships have a price that moves.
\iqlookfor{the two break-evens and the capital at risk in a purchase.}
\end{solution}

\begin{solution}{iq:m1:getting-access-futures-options:5}
Negotiable: commission per side and its minimum; the margin multiple and
intraday limits; interest on excess cash and the collateral accepted;
\omterm{def:m1:getting-access-equities:brokers}{give-up} and platform fees; notice periods. Not negotiable: exchange and
regulatory fees, the clearing house's margin, segregation rules. The firm
brings volume, balances, a clean risk profile and an alternative FCM.
\iqlookfor{margin terms ahead of commission.}
\end{solution}

\begin{solution}{iq:m1:getting-access-futures-options:6}
Fixed costs and scale. A futures firm needs one membership, one feed, one
book. An \omterm{def:m1:options-market-structure:mm}{options market maker} needs appointments and permits on a dozen or
more exchanges for the same series, capacity to quote hundreds of thousands
of series and consume a feed of hundreds of billions of messages, models
for the whole surface, capital for inventory that cannot be flattened at the
close, and enough breadth of flow for the inventory to net. Entitlements
reward incumbents. The minimum efficient scale is high, so the industry is
concentrated.
\iqlookfor{inventory and breadth, not just technology.}
\end{solution}
